\chapter{LITERATURE REVIEW AND EXISTING SYSTEMS}
\thispagestyle{empty}

This chapter reviews existing approaches in learning management, OCR
document conversion, AI-assisted content generation, deterministic
extraction, and multi-tenant architectures, and identifies the gap addressed
by this project.

\clearpage

\section{Existing Approaches in the Domain}

\subsection{Learning management and assessment platforms}

Commercial and open-source learning management platforms provide course
organization, content delivery, and basic assessment features. They are
generally oriented toward enrolled students and instructor-authored quizzes;
connecting arbitrary printed study material, OCR conversion, AI-assisted
content generation, and examination-paper assembly into a single pipeline is
not their primary purpose. As a result, institutions that need the full
content-to-examination workflow typically combine several tools and
manually move data between them.

\subsection{OCR for document conversion}

Optical character recognition converts scanned documents and images into
machine-readable text. General-purpose OCR engines trade recognition accuracy
against resource cost and must be tuned for low-resource deployments.
PaddleOCR (PP-OCR) addresses this with an ultra-lightweight pipeline designed
to run on limited hardware \cite{ppocr,ppocrv3}; it supports per-page and
whole-document recognition with a detection--recognition model pair. A
practical engineering strategy is to extract embedded text directly from
PDF files first---using a library such as PyMuPDF \cite{pymupdf}---and to
fall back to full OCR only for pages that lack extractable text, which
avoids the cost of re-OCR of born-digital documents.

\subsection{Large-language-model-assisted content generation}

Large language models can draft prose, questions, summaries, and
outlines, but their open-ended output is unreliable for structured academic
assets unless constrained. The reliable pattern used in practice is to
require strictly schema-validated output, validate it at every trust boundary
(the worker and the API), request a re-generation on validation failure, and
keep a human reviewer in the loop for anything that becomes authoritative.
Structured-data validation libraries in both the JavaScript ecosystem (Zod)
and the Python ecosystem (Pydantic) \cite{pydantic} provide the canonical
schemas for this enforcement.

\subsection{Deterministic extraction versus machine learning}

For tasks such as recovering the structure of an examination paper or
extracting question candidates from text, opaque machine-learning predictions
complicate review: an institution needs an explicit report of what was
recognized and where the extractor is uncertain. Deterministic, rule-based
extractors that resolve ambiguity explicitly are easier to audit and correct,
and they pair well with AI where AI supplies suggestions that the teacher
confirms.

\subsection{Multi-tenant SaaS architectures}

Multi-tenant systems serve many customers from a single deployment while
isolating their data. A modular monolith---one deployable containing clearly
bounded modules---avoids premature microservice fragmentation while still
permitting independently deployable components where scaling genuinely
requires it. In this project the API is the modular monolith and the OCR
service is the single separately deployable component.

\section{Limitations of Existing Approaches}

Existing systems address fragments of the pipeline but not its integration.
Common limitations include: manual movement of data between disconnected
tools; OCR pipelines that are centralized, expensive to scale, and hard to
retry per page; AI output that is not schema-validated and therefore cannot
be safely imported into structured domains; grading and extraction that is
opaque or non-deterministic; and little or no support for strict tenant
isolation when several institutes share a deployment.

\section{Gap Addressed by the Project}

The gap addressed by CatLium EduTech is an integrated, deterministic,
tenant-isolated pipeline that runs from syllabus and raw material to a
delivered and evaluated examination. The design choices address the
limitations above directly: OCR is a separately deployable service with
per-page fallback and a distributed worker model; AI output is validated on
both worker and API boundaries and always requires human review; extraction
and grading are deterministic and produce explicit reports; and every data
access is tenant-scoped.